\documentclass[11pt]{article}

\usepackage[utf8]{inputenc}
\usepackage[T1]{fontenc}
\usepackage[margin=1in]{geometry}
\usepackage{booktabs}       % professional-quality tables
\usepackage{tabularx}       % tables that fill the text width
\usepackage{hyperref}       % hyperlinks
\usepackage{xcolor}         % colors
\usepackage{microtype}      % microtypography
\usepackage{parskip}        % blank line between paragraphs, no indent
\usepackage{enumitem}       % list spacing control

\hypersetup{
  colorlinks = true,
  linkcolor  = black,
  urlcolor   = blue!55!black,
}

% Left-aligned ragged-right column that fills remaining width.
\newcolumntype{L}{>{\raggedright\arraybackslash}X}

\title{Lifelong Assistant}
\author{Swanand Pande}
\date{Drafted 28 August 2026 \\ Revised 10 September 2026}

\begin{document}

\maketitle


\section{Introduction}

A person has some preferences about activities and things around. And it is not necessary that preferences remain static. These preferences can change over time. A person liking espresso today might prefer a latte a month later. 

As people start using assistants to get help in their everyday tasks, it would be important for these assistants to capture the evolving nature of preferences. This would give a better UX to the user.

The project is to create an assistant that knows facts and preferences of the user. Importantly, the assistant knows the evolving nature of preferences. 

\section{Scope}

\begin{enumerate}[itemsep=2pt, topsep=4pt]
  \item \textbf{Dataset 1} --- assistant/user conversations containing facts and
        evolving preference patterns.
  \item \textbf{Dataset 2} --- a learning-trace dataset, in which the assistant evolves
        its understanding one day at a time.
  \item Generate a response to a new user message.
  \item Initiate a new conversation.
\end{enumerate}

\section{Current state}

\subsection{Ground truth}

\begin{table}[h]
\centering
\begin{tabularx}{\textwidth}{@{}l L@{}}
\toprule
\textbf{Entity} & \textbf{Rule} \\
\midrule
\texttt{coffee}   & Fixed 7-day cycle \\
\texttt{clothing} & Alternates every 3 days \\
\texttt{exercise} & Day-of-week: Mon/Wed/Fri $\to$ morning yoga,
                    Tue/Thu $\to$ evening run \\
\bottomrule
\end{tabularx}
\end{table}

\subsection{Main evaluation dataset}

\begin{table}[h]
\centering
\begin{tabularx}{\textwidth}{@{}l L@{}}
\toprule
Generator      & gpt-4o-mini \\
Sessions       & 60 (1 session = 1 day) \\
QA pairs       & 67 \\
Question types & factual, factual\_evolving, recall, pattern\_identification,
                 prediction, transition \\
\bottomrule
\end{tabularx}
\end{table}

\section{Things to do}

\subsection{(a) Agentic approach}

Do not hardcode the context. Expose the tools and let the agent decide what to call.

The prior design ran a single fixed pass per query. That pass does not guarantee the
result is adequate. The agentic version replaces the pass with a loop.

At each step, the agent asks two questions:

\begin{enumerate}[itemsep=2pt, topsep=4pt]
  \item Is this new, or does it update something already known?
  \item Is the retrieved context enough to answer the query?
\end{enumerate}

\subsubsection*{Tools}

\begin{table}[h]
\centering
\begin{tabularx}{\textwidth}{@{}l l L@{}}
\toprule
\textbf{Tool} & \textbf{Mode} & \textbf{Purpose} \\
\midrule
\texttt{list\_entities}        & both  & Read the live taxonomy before writing or querying \\
\texttt{list\_preferences}     & write & Inspect what is already recorded \\
\texttt{add\_preference}       & write & Record one dated choice \\
\texttt{remove\_preference}    & write & Correct a duplicate or mis-filed entry \\
\texttt{index\_conversation}   & write & Add a chunk to the RAG store \\
\texttt{summarize\_week}       & write & Per-entity transitions plus facts \\
\texttt{summarize\_month}      & write & Consolidate weeks into candidate rules \\
\texttt{summarize\_year}       & write & Confirm which rules held \\
\texttt{summarize\_lifetime}   & write & Standing profile \\
\texttt{build\_context}        & read  & One-call assembly for a query \\
\texttt{search\_preferences}   & read  & Exact dated timeline \\
\texttt{search\_conversations} & read  & Verbatim retrieval \\
\texttt{get\_summary}          & read  & Fetch one summary by level \\
\bottomrule
\end{tabularx}
\end{table}

Summarization is deliberately split by level rather than exposed as one call, so the
agent chooses how far up the hierarchy to roll and the evaluation can attribute cost to
a specific level.

Context assembly holds no domain knowledge: it reads the entity taxonomy out of the
store at call time and matches the query against whatever is actually there, then
discards any entity name the model returns that does not literally exist. The earlier
builder carried a hand-written table of expected domains with per-domain keyword lists
--- which works only on a dataset whose categories are known in advance, and silently
returns nothing for anything else.

\subsection{(b) Continual learning}

The agent should build its understanding on the go --- one day at a time, using only
what it has seen so far.

\end{document}
